\section*{Chapitre \ref{ch:rc:the-valuation-adjustment-desk} --- Le desk d'ajustements de valeur}

\begin{solution}{exo:rc:the-valuation-adjustment-desk:1}
Compensation~: les valeurs de la nouvelle transaction compensent celles de l'ensemble existant sur de nombreux chemins, si bien que l'exposition positive nette baisse.
Isolée, la transaction a sa propre exposition positive~; dans l'ensemble, elle retranche
davantage d'exposition qu'elle n'en ajoute.
\end{solution}

\begin{solution}{exo:rc:the-valuation-adjustment-desk:2}
$1\,000\,000/(100\,000\,000\times13.67) = 7.3$ points de base par an.
\end{solution}

\begin{solution}{exo:rc:the-valuation-adjustment-desk:3}
Une mesure de la qualité de crédit (telle qu'une notation), du secteur d'activité et de la région géographique.
\end{solution}

\begin{solution}{exo:rc:the-valuation-adjustment-desk:4}
$\mathrm{CVA} = (1-R)\sum_k\E[D\max(\sum_iV_i,0)]\Delta\mathrm{PD}_k$ et
$\max(\sum_iV_i,0) = \sum_iV_i\mathbf 1_{\sum_jV_j>0}$ chemin par chemin~; on prend les espérances et on somme sur $i$~:
les contributions $a_i$ s'additionnent. De façon équivalente, la CVA est homogène de degré un par rapport aux tailles, et
le théorème d'Euler s'applique.
\end{solution}

\begin{solution}{exo:rc:the-valuation-adjustment-desk:5}
La charge incrémentale d'une transaction dépend de ce qui figure déjà dans l'ensemble, si bien que deux transactions opposées
sont facturées différemment selon celle qui est arrivée en premier. Les contributions d'Euler se calculent sur
l'ensemble entier, d'un seul coup, et sont symétriques entre les transactions, de sorte qu'elles ne dépendent pas de l'historique.
\end{solution}

\begin{solution}{exo:rc:the-valuation-adjustment-desk:6}
La base entre le spread du client et l'indice, le \omterm{def:rc:reduced-form-credit:jtd}{risque de saut au défaut} du client (la protection
sur indice paie peu lorsqu'un seul nom fait défaut), le \omterm{def:rc:counterparty-exposure:wwr}{risque de corrélation défavorable}, et les variations de l'exposition elle-même.
\end{solution}

\begin{solution}{exo:rc:the-valuation-adjustment-desk:7}
2\,385~USD par point de base~; divisé par l'\omterm{def:rc:reduced-form-credit:risky}{annuité risquée} à dix ans de l'émetteur comparable et par un point de base,
soit environ 3.55~millions de dollars de protection.
\end{solution}

\begin{solution}{exo:rc:the-valuation-adjustment-desk:8}
La charge est la valeur de l'ajustement à l'origine~; ensuite, le P\&L du desk évolue avec le
spread du client, les facteurs de l'exposition, le financement et les coûts en capital, ainsi que tout écart non couvert. Seule une
couverture parfaite de chaque composante la laisserait inchangée, et le financement comme le capital ne peuvent pas être couverts.
\end{solution}

\begin{solution}{pb:rc:the-valuation-adjustment-desk:1}
\textbf{1.} 4.11~\%~; une annuité de 13.67.
\textbf{2.} 337.5 points de base, issus de la régression sur émetteurs comparables, le client étant rattaché à BB,
secteur industriel, États-Unis.
\textbf{3.} CVA 1\,649\,494~USD (12.1 points de base), FCA 343\,872 (2.5), KVA 3\,141\,826 (23.0).
\textbf{4.} 37.6 points de base par an.
\textbf{5.} KVA~: un client non noté reçoit la pondération de risque CVA de 7~\% et un swap à vingt ans immobilise
du capital pendant des décennies.
\textbf{6.} La CVA, soit 32~\% du total~; le financement et le capital ne peuvent pas être couverts sur le marché.
\textbf{7.} 2\,385~USD par point de base~; environ 3.55~millions de dollars de protection sur l'émetteur comparable à dix ans.
\textbf{8.} La base avec le crédit propre du client, son \omterm{def:rc:reduced-form-credit:jtd}{risque de saut au défaut}, la croissance de l'exposition au-delà
dix ans, et le \omterm{def:rc:rates-risk:crossgamma}{gamma croisé}.
\textbf{9.} Rarement pour une entreprise non notée~; si elle était disponible, elle coûterait environ la CVA plus la
marge du vendeur et le capital.
\textbf{10.} Avec des swaps et des swaptions dimensionnés sur le delta de taux de la CVA et son vega, sur des
scénarios communs.
\textbf{11.} Une meilleure notation interne, un spread de financement plus bas, un \omterm{def:rc:funding-margin-and-capital-adjustments:kva}{taux de rendement minimal exigé} plus bas, un
autre traitement du capital CVA, ou la tarification du seul effet incrémental sur une grande
relation existante.
\textbf{12.} Avec BBB et la pondération de risque de 3~\%, la surcharge tombe à 26.2 points de base (CVA 7.2,
FCA 3.1, KVA 15.8).
\textbf{13.} Un CSA, une clause de résiliation anticipée à dix ans, un swap plus court renouvelé, ou du collatéral.
\textbf{14.} Elle tarifie peut-être mal le crédit du client ou ne facture pas le capital, et elle
attirera les clients dont elle sous-tarife le risque.
\textbf{15.} Seulement si la transaction aurait rémunéré ses risques~; renoncer à une
transaction sous-tarifée évite une perte.
\textbf{16.} La banque fixe le \omterm{def:rc:funding-margin-and-capital-adjustments:kva}{taux de rendement minimal exigé}~; le \omterm{def:rc:the-valuation-adjustment-desk:desk}{desk XVA} devrait le facturer au desk d'origine afin que
la rentabilité de la transaction soit visible là où la décision se prend.
\textbf{17.} Le crédit de vente devrait être diminué de la \omterm{def:rc:the-valuation-adjustment-desk:charge}{charge XVA}, afin que le commercial ne soit pas rémunéré pour
un revenu que les risques de la banque consomment.
\textbf{18.} Que le taux intègre le coût de son crédit, celui du financement d'une position non garantie et
du capital que la banque détient, et qu'un CSA ou une clause de résiliation anticipée le ferait baisser.
\textbf{19.} Résultat à retenir~: \emph{l'entreprise non notée}~: une surcharge globale de 37.6 points de base par
an, dont le desk peut couvrir sur le marché les 12.1 points de base de CVA (32~\%).
\textbf{20.} La gestion des coûts de contrepartie, de financement et de capital que nul desk isolé ne peut
voir ni couvrir seul.
\end{solution}

\begin{solution}{iq:rc:the-valuation-adjustment-desk:1}
Les ajustements sont des propriétés d'\omterm{def:rc:counterparty-exposure:netting}{ensembles de compensation} qui traversent plusieurs desks~; seule une vue centrale peut valoriser
la compensation et le collatéral, éviter le double comptage, couvrir le crédit de façon centralisée et gérer le financement et
le capital de manière cohérente.
\iqlookfor{\omterm{def:rc:counterparty-exposure:netting}{ensembles de compensation} entre desks et couverture centralisée.}
\end{solution}

\begin{solution}{iq:rc:the-valuation-adjustment-desk:2}
La contribution de chaque transaction est la dérivée du total par rapport à un changement d'échelle de cette transaction~;
pour une fonction homogène de degré un, les contributions s'additionnent au total (théorème d'Euler).
La CVA à crédit fixe est homogène~; les formules de capital avec planchers ou plafonds ne le sont peut-être pas.
\iqlookfor{la définition de la dérivée et la condition d'homogénéité.}
\end{solution}

\begin{solution}{iq:rc:the-valuation-adjustment-desk:3}
Construire une courbe de \omterm{def:rc:the-valuation-adjustment-desk:proxy}{spreads de crédit de substitution} à partir d'émetteurs liquides comparables par notation, secteur et région (régression ou classes), ou
la rattacher à un nom apparenté~; valider sur des obligations ou des prêts s'il y en a~; couvrir avec une protection sur indice ou
sur émetteur comparable et conserver des réserves au titre de la base.
\iqlookfor{construction des spreads de substitution et base résiduelle.}
\end{solution}

\begin{solution}{iq:rc:the-valuation-adjustment-desk:4}
Conserver en mémoire, pour chaque \omterm{def:rc:counterparty-exposure:netting}{ensemble de compensation}, les valeurs des chemins sur les scénarios communs~; valoriser la nouvelle transaction sur
les mêmes chemins (formules fermées ou grilles précalculées)~; recalculer les profils nets et
les ajustements de façon incrémentale~; paralléliser~; se replier sur des approximations pour les transactions exotiques.
\iqlookfor{chemins stockés et recalcul incrémental.}
\end{solution}

\begin{solution}{iq:rc:the-valuation-adjustment-desk:5}
Les variations des spreads de contrepartie et des expositions, nettes des couvertures, des coûts de financement, du capital, des charges liées aux
nouvelles transactions et des écarts de couverture. En situation de crise, spreads et expositions bondissent ensemble (\omterm{def:rc:counterparty-exposure:wwr}{risque de corrélation défavorable}),
les spreads de substitution se découplent des clients, les couvertures deviennent coûteuses et les couvertures du capital CVA peuvent elles-mêmes
être facturées.
\iqlookfor{les moteurs du P\&L et la dynamique en période de crise.}
\end{solution}

\begin{solution}{iq:rc:the-valuation-adjustment-desk:6}
En totalité, la banque perd les clients sensibles au prix au profit de banques aux hypothèses plus favorables et conserve
ceux dont le risque est tarifé plus haut par d'autres (sélection adverse)~; en partie, elle subventionne des transactions. Les banques
répercutent l'essentiel de cet effet, se font concurrence sur les hypothèses et la valeur des relations, et doivent être cohérentes.
\iqlookfor{concurrence, sélection adverse et cohérence.}
\end{solution}
